\section{Introduction}

\textit{Web browsers} are central to modern computing, with
browser-based usage corresponding to over \textbf{80\%} of daily work
in some organizations~\cite{montenegro2025workforce}. Browser makers
therefore invest heavily in fuzzing, audits, security research,
standardization, shared test suites, and interoperability efforts aimed
at ensuring browser correctness~\cite{google-project-zero,mozilla-mythos,w3c,whatwg,wpt,interop}.

Developers increasingly use \textit{fuzzing}, an automated testing
technique that exercises software with large numbers of randomly
generated or mutated inputs, to identify and fix bugs. Browser fuzzers
have been successful at finding crashes, security vulnerabilities, and
rendering regressions. R2Z2 detects rendering regressions by fuzzing the
output of different browser versions~\cite{song2022r2z2}; Metamong
detects render-update bugs through fuzzing~\cite{song2023metamong}; and
FuzzOrigin targets UXSS vulnerabilities through origin
fuzzing~\cite{kim2022fuzzorigin}. However, one important browser
mechanism remains difficult to test systematically: \textit{layout
invalidation}, which determines which parts of a page must be
recomputed after changes to the DOM or CSS.

Layout invalidation is difficult because browsers attempt to reuse as
much previous work as possible while still ensuring that the final page
state is correct. The browser rendering pipeline includes parsing,
style computation, layout, painting, and related stages~\cite{wbe}.
After a DOM or CSS update, invalidation decides which work can be reused
and which work must be recomputed. Public browser-engine trackers
contain invalidation failures in Chromium, Firefox, and WebKit-based
browsers~\cite{chromium497183437,chromium379886422,mozilla1452426,mozilla539356,webkit233421,webkit97801}.
Modern layout engines, including Chrome's LayoutNG, are designed in part
to reduce these correctness bugs~\cite{kilpatrick2021layoutng}, and
some optimizations have historically stalled because ensuring correct
invalidation was too difficult~\cite{mozilla2006reflow}.

This paper introduces \toolname, a fuzzer for detecting
under-invalidation bugs in modern web browser engines.
\toolname systematically exercises incremental layout behavior
by generating DOM and CSS updates that should force the browser to
recompute affected rendering state. Its central strategy is
self-comparison: the system compares the browser's incremental layout
result against its own from-scratch layout result for the same final
page. This allows \toolname to detect stale layout state without
relying on a separate browser implementation.

This paper makes the following contributions:
\begin{itemize}
    \item We describe the challenges of extending fuzzing to browser
    layout invalidation and explain why existing browser fuzzers do not
    systematically target under-invalidation bugs.
    \item We design \toolname, a targeted fuzzer that generates
    DOM and CSS updates, compares incremental and fresh layout results,
    and reports stale layout or rendering state.
    \item We evaluate \toolname through experiments and ablation
    studies across XXX browser engines, measuring bug discovery,
    minimization, and clustering effectiveness.
    \item We release \toolname together with evaluation
    artifacts, generated test cases, and benchmarks at the following
    URL: XXX.
\end{itemize}
